\documentclass{article}
\usepackage{amsmath}
\usepackage{cancel}


\title{Velocity-Dependent Potentials}

\begin{document}
Previously, we considered forces that are derivable from a potential function,
that is
\begin{equation}
  \vec{F}_i^{(a)} = -\nabla_iV
\end{equation}
We can actually do better than this and generalize our potential function. What
if the potential depends not only on the coordinates $\{q_j\}$ but also the
coordinate velocities $\{\dot{q}_j\}$? In this case, define
\begin{equation}
  Q_j = -\frac{\partial U}{\partial q_j}+\frac{d}{dt}\left( \frac{\partial U}{\partial \dot{q}_j} \right).
\end{equation}
D'Alembert's principle says
\begin{equation}
  \frac{d}{dt}\left( \frac{\partial T}{\partial \dot{q}_j} \right)-\frac{\partial T}{\partial q_j}-Q_j = 0
\end{equation}
From this, we can see that our generalized potential still allows us to write
\begin{equation}
  \mathcal{L}= T-U
\end{equation}
so that d'Alembert's principle again reduces to the Euler-Lagrange equations:
\begin{equation}
  \frac{d}{dt}\frac{\partial \mathcal{L}}{\partial \dot{q}_j}-\frac{\partial \mathcal{L}}{\partial q_j} = 0 
\end{equation}
\end{document}
